\section*{Discussion}

NeuroscienceBench provides a quantitative measure of how reliably AI agents perform the analytical and interpretive decisions that are common in neuroscience research. Currently, the strongest model configurations are Opus 5.5 and Opus 5 under Pi that passed 32.2\% and 30.8\% of attempts. Harness contributed less than model identity, yet the same model shifted by several points between harnesses. Although aggregate pass rates were similar across most categories, model rankings were not stable. The strongest category-dependent effect was data type: mean pass rates ranged from 14\% on less common omics readouts and 18\% on clinical and biofluid data to about 25\% on behavior and on EEG and sleep scoring, and individual models swung by about 28\% between their best and worst data type.

These results have several limitations. The evaluations are retrospective and built from published data; anti-shortcut checks reduce, but cannot exclude, the possibility that a model has encountered the source study. With 91 evaluations distributed over twelve functional systems and twelve process classes, many categories contain fewer than ten evaluations, and each configuration made only three attempts, so category-level estimates can be imprecise and should be treated as hypotheses for targeted expansion. Finally, initial difficulty was calibrated on two frontier models, Claude Opus 5.5 and GPT-6 Astra. The two models nevertheless fared very differently. Opus 5.5 ranked first overall, whereas GPT-6 Astra ranked eleventh under Pi and eighteenth under Codex.
